\section{Overview}

The analog front-end site (\peripheral{AFEx}) is the digital half of one electrochemical channel. Four identical sites, one per channel tile, drive the four \texttt{anatop\_pixel} cells of the analog macro: a grounded-working-electrode potentiostat (control amplifier into the counter electrode, transimpedance amplifier holding the working electrode at the common mode), a 320-code programmable feedback ladder, a 12-bit potential-program DAC, a 12-bit common-mode DAC, a local bias generator, a 16:1 converter input multiplexer, a 16:1 analog test-port multiplexer and a 10-bit SAR converter. The site owns the 50 control bits of that cell and sequences its converter. The main features are listed below:

\begin{itemize}
	\item The 50 control bits as direct register outputs: \register{AFExTIA} (ladder code), \register{AFExDACVP} and \register{AFExDACVCM} (DAC codes and enables), \register{AFExBIAS} (bias trim), \register{AFExMUX} (converter and test-port selects)
	\item A converter sequencer: trigger clock at MCLK\,/\,$2(\register{AFECLKDIV}+1)$, twelve falling edges per conversion in $16 + \register{AFESAMPLESTEP}$ ticks, reset release with a 128-cycle quiet period
	\item Result capture inside the READY window with the converter's inverted bit 9 corrected, into a four-entry FIFO tagged with the input slot and the excitation phase (\register{AFExDATA})
	\item Single-shot and continuous conversion (\register{AFESTART}, \register{AFECONT}), a simultaneous-sample trigger shared by the four sites (\register{AFESYNC}, \register{AFESYNCEN}), and a two-code excitation swap engine (\register{AFExSWAP}, \register{AFESWAPEN})
	\item One level interrupt per site, combined onto vector 124 (\register{AFEDRDYIE}, \register{AFEERRIE})
	\item Per-site ownership: a site answers its channel hart or hart 0; every other master reads 0 and its writes are dropped
\end{itemize}

\section{Functional description}

\subsection{Control surface}

Every field of \register{AFExTIA}, \register{AFExDACVP}, \register{AFExDACVCM}, \register{AFExBIAS} and \register{AFExMUX} reaches the analog macro directly as 1.0\,V CMOS; the 1.0\,V to 2.5\,V level shifters are inside the macro. The bits reset to zero, which is the minimum-gain, converter-on-the-transimpedance-output state, and two of the reset values are not usable operating points: \register{AFEBIASADJ} = 0 is not the nominal bias trim, and \register{AFEATPSEL} = 0 selects an undriven test-port slot. The power-up sequence in the Programming section loads them before any measurement.

The feedback resistance follows $R_f = 18\,\mathrm{k\Omega} + 6\,\mathrm{k\Omega}\,N$ with $N = \register{AFERESEN} + 64\,\mathrm{popcount}(\register{AFETHEN})$, $N = 0 \ldots 319$; the law is 9\,\% high at $N = 0$ and within 0.5\,\% from $N = 20$ up, so gain constants are stored per code, never computed. Both DACs step 610.35\,\textmu V per code.

\subsection{Converter sequencer}

\register{AFEEN} releases the converter reset (SARADC\_rst = not EN) and starts a 128-cycle quiet period; \register{AFEREL} reports its end and \register{AFESTART} is ignored before it. A conversion is $16 + \register{AFESAMPLESTEP}$ ticks of the trigger clock: one idle tick, the clear pulse (one tick), one low tick, the sample pulse ($1 + \register{AFESAMPLESTEP}$ ticks; its falling edge is the aperture), one low tick, ten conversion pulses (one falling edge per bit trial) and one low tick. With \register{AFESAMPLESTEP} = 7 and a 20\,MHz trigger clock that is 23 ticks, 1.150\,\textmu s, 869.6\,kSa/s, the pattern the macro is characterised with. The trigger clock is MCLK\,/\,$2(\register{AFECLKDIV}+1)$: a 40\,MHz MCLK at \register{AFECLKDIV} = 0 gives 20\,MHz; a 24\,MHz MCLK gives 12\,MHz, which is slower than characterised and inside the macro's synchronous margin. The output is a register, so the gated clock is glitch-free.

The macro raises READY after the tenth conversion falling edge and holds the data bus valid for about 100\,ns. The site captures the bus on the MCLK edge after the one that first samples READY high (25 to 50\,ns after the rise at 40\,MHz), inverts bit 9, and pushes the code into the FIFO with the \register{AFEADCSEL} value and the swap phase latched at the aperture. A READY missing for 2048\,MCLK after an aperture, or a new aperture with the previous result still outstanding, sets \register{AFETO}; a capture into a full FIFO sets \register{AFEOVF} and drops the sample.

\subsection{Result FIFO}

\register{AFExDATA} is the FIFO head and reading it has no side effect: \register{AFEVALID} says whether it holds a result and \register{AFECNT} counts the entries. Writing 1 to \register{AFEDRDY} pops the head; a pop on an empty FIFO does nothing. The tags let firmware interleave the current channel (\register{AFEADCSEL} = 0) and the potential channel (\register{AFEADCSEL} = 2) in one continuous run and sort the results afterwards.

\subsection{Simultaneous sampling and the swap engine}

A write of 1 to \register{AFESYNC} on any site pulses a trigger shared by the four sites, and every site with \register{AFESYNCEN} set starts on the same MCLK edge; sites programmed with the same \register{AFECLKDIV} then run tick-aligned. With \register{AFESWAPEN} and \register{AFECONT} set, A\_Dac\_Vp alternates between \register{AFEVP} and \register{AFEVP2} every $\register{AFEPERIOD} + 1$ conversions, switching at a conversion boundary, so the square-wave excitation and the converter strobe share one divider and the sampling grid is commensurate with the excitation. The phase at each aperture is recorded in \register{AFEPHTAG}.

\section{Interrupts and events}

Each site drives one level interrupt, $(\register{AFEDRDY} \wedge \register{AFEDRDYIE}) \vee ((\register{AFEOVF} \vee \register{AFETO}) \wedge \register{AFEERRIE})$, and the four lines are ORed onto vector 124 (the interrupt router caps the source count at 127, which leaves no room for four). Every channel hart routes vector 124; the service routine reads its own site's \register{AFExSR}, returns if nothing is set there, otherwise pops the FIFO until \register{AFECNT} is 0 and clears \register{AFEOVF} and \register{AFETO} with a write of 1. The block produces no event-fabric taps.

\section{Programming}

\subsection{Power-up}

\begin{enumerate}
	\item Write the stored bias trim to \register{AFEBIASADJ}; nothing else on the channel is meaningful before the bias rails are up.
	\item Enable and program the common-mode DAC (\register{AFEVCMEN}, \register{AFEVCM}) and the potential DAC (\register{AFEVPEN}, \register{AFEVP}).
	\item Set the gain code in \register{AFExTIA} and \register{AFEADCSEL} = 0.
	\item Set \register{AFEEN} and poll \register{AFEREL}.
\end{enumerate}

\subsection{One conversion}

Write \register{AFESTART}; wait for \register{AFEDRDY} (or the interrupt with \register{AFEDRDYIE}); read \register{AFExDATA}; write 1 to \register{AFEDRDY} to pop. Wait 0.54\,\textmu s after any \register{AFEADCSEL} change before the next start.

\subsection{Continuous record}

Set \register{AFECONT} with \register{AFESTART} (or arm \register{AFESYNCEN} on every site and write \register{AFESYNC} once for a simultaneous record). Drain the FIFO faster than one conversion per 1.15\,\textmu s at the design rate, or set \register{AFEERRIE} and treat \register{AFEOVF} as a lost sample. Clear \register{AFECONT} to stop at the next conversion boundary; clear \register{AFEEN} to abort and reset the converter.

\subsection{Square-wave impedance record}

Program \register{AFEVP} and \register{AFEVP2} as a code pair inside one octave, \register{AFEPERIOD} for the excitation half period in conversions, then \register{AFESWAPEN} with \register{AFECONT} and \register{AFESTART}. Discard the electrode's own transient (at least 100 time constants), then correlate the tagged results over an integer number of periods.

\section{Usage notes}

\begin{itemize}
	\item Reads of every register are side-effect-free; the FIFO is popped by \register{AFEDRDY} only, so LR/SC and AMO sequences to a site address are harmless.
	\item \register{AFESTART} and \register{AFESYNC} read 0. A \register{AFESTART} written in the same access that sets \register{AFEEN} is lost, because \register{AFEREL} is not yet set.
	\item The converter's raw bus carries bit 9 inverted; \register{AFECODE} is already the straight code, offset-binary about the common-mode code. Do not invert again.
	\item Converter input slots 3 and 6 to 15 are undriven in the tapeout cell; select only 0, 1, 2, 4 and 5.
	\item The data hold time of the converter bus after READY is not characterised; the capture point assumes at least two MCLK periods and is to be confirmed on the extracted macro before tapeout.
\end{itemize}
